% Part: first-order-logic
% Chapter: tableaux
% Section: soundness-identity

\documentclass[../../../include/open-logic-section]{subfiles}

\begin{document}

\olfileid{fol}{tab}{sid}

\olsection{\usetoken{S}{identity} સાથે યથાર્થતા}

\begin{prop}
તાદાત્મ્ય માટેના નિયમો ધરાવતા !!^{tableau}s યથાર્થ છે: કોઈ બંધ
!!{tableau} સંતોષ્ય નથી.
\end{prop}

\begin{proof}
પહેલાંની જેમ આપણે માત્ર એટલું દર્શાવવું છે કે જો !!a{tableau}માં
સંતોષ્ય શાખા હોય, તો તેને~$\eq$ માટેના નિયમોમાંથી કોઈ એક લાગુ કરતાં
મળતી શાખા પણ સંતોષ્ય છે. શાખા પરનાં ચિહ્નિત !!{formula}sનો ગણ
$\Gamma$ હોય અને~$\Gamma$ને સંતોષતી !!a{structure} $\Struct{M}$ હોય.

ધારો કે શાખાને~$\eq$ વાપરીને, એટલે કે ચિહ્નિત
!!{formula}~$\sFmla{\True}{\eq[t][t]}$ ઉમેરીને, વિસ્તારીએ છીએ. દેખીતી
રીતે $\Sat{M}{\eq[t][t]}$, તેથી $\Struct{M}$ એ
$\Gamma \cup \{\sFmla{\True}{\eq[t][t]}\}$ને પણ સંતોષે છે.

જો શાખાને $\TRule{\True}{\eq}$ વડે વિસ્તારીએ, તો ચિહ્નિત
!!{formula}~$\sFmla{\True}{!A(t_2)}$ ઉમેરીએ છીએ, પરંતુ $\Gamma$માં
$\sFmla{\True}{\eq[t_1][t_2]}$ અને $\sFmla{\True}{!A(t_1)}$ બંને છે.
તેથી $\Sat{M}{\eq[t_1][t_2]}$ અને $\Sat{M}{!A(t_1)}$. એવી ચલ-સોંપણી
$s$ લો કે $s(x) = \Value{t_1}{M}$.
\olref[syn][ass]{prop:sentence-sat-true} મુજબ
$\Sat{M}{!A(t_1)}[s]$. $\varAssign{s}{s}{x}$ હોવાથી
\olref[syn][ext]{prop:ext-formulas} મુજબ $\Sat{M}{!A(x)}[s]$.
$\Sat{M}{\eq[t_1][t_2]}$ હોવાથી
$\Value{t_1}{M} = \Value{t_2}{M}$ અને તેથી
$s(x) = \Value{t_2}{M}$. ફરીથી
\olref[syn][ext]{prop:ext-formulas} લાગુ કરતાં
$\Sat{M}{!A(t_2)}[s]$ પણ મળે છે.
\olref[syn][ass]{prop:sentence-sat-true} મુજબ $\Sat{M}{!A(t_2)}$.
$\TRule{\False}{\eq}$નો કિસ્સો એ જ રીતે સંભાળીએ છીએ.
\sourcecorrection{OLTAB-012}{મૂળની \texttt{soundness-identity.tex},
પંક્તિ~31માં $\TRule{\True}{\eq}$ના ફલિતનું ચિહ્ન અનિર્ધારિત~$S$
લખાયું છે. પ્રદર્શિત નિયમ અને આ જ દલીલનાં બંને આધારવાક્યો મુજબ અહીં
$\True$ પુનઃસ્થાપિત કર્યું છે.}
\end{proof}

\end{document}
